\documentclass[../main.tex]{subfiles}
\graphicspath{{\subfix{../images/}}}

\begin{document}
    \section{Análisis de \texorpdfstring{\emph{L\textsubscript{2}}}{L2}}
    \[
        \boxed{L_2=\{a^{2n}b^3:n\ge0\}}
    \]

    El prefijo ``a'' crece de dos en dos respecto a \(n\). Si analizamos las cadenas generadas por \(n\in\{0,1,2\}\) obtenemos:
    \begin{align*}
        w_0 &= \text{``bbb''}\\
        w_1 &=\text{``aabbb''}\\
        w_2 &=\text{``aaaabbb''}
    \end{align*}

    Entonces, cualquier cadena \(w \in L_2\) comenzará con una cantidad par de aes y terminará siempre con exactamente tres bes. La máquina de Turing \(M_2\) verifica que se cumplan estas condiciones avanzando hasta el símbolo blanco (\textcrb) que representa el final de la cadena, validando el sufijo ``bbb'' y consumiendo en bucle letras aes de izquierda a derecha hasta encontrar el blanco inicial. \(M_2\) alcanzará su estado de aceptación únicamente si la cantidad de aes encontrada es par.
\end{document}
